
\setcounter{chapter}{1}
\chapter{Design}
\minitoc %insert la minitoc
\graphicspath{{Chapitre2/figures/}}

%\DoPToC

%==============================================================================
\pagestyle{fancy}
\fancyhf{}
\fancyhead[R]{\bfseries\rightmark}
\fancyfoot[R]{\thepage}
\renewcommand{\headrulewidth}{0.5pt}
\renewcommand{\footrulewidth}{0pt}
\renewcommand{\chaptermark}[1]{\markboth{\MakeUppercase{\chaptername~\thechapter. #1 }}{}}
\renewcommand{\sectionmark}[1]{\markright{\thechapter.\thesection~ #1}}

\begin{spacing}{1.2}
%==============================================================================
\section*{Introduction}
La partie conception de l'application \textbf{\textsl{n'est pas toujours obligatoire}}. En effet, quand notre travail consiste en une étude théorique, ou une mise en place d'un système par exemple, 
il est inutile voire obsolète de faire un diagramme de classes ou de séquence.\\
Quand il s'agit de développement, par contre, la partie conception s'impose. 
\section{Recommadations}
En général, 
il faut suivre les règles suivantes :
\begin{itemize}
	\item Choisir une méthodologie de travail : un processus unifié, une méthode agile;
\end{itemize}
\section{Diagrammes}
Il faut Bien choisir les diagrammes adéquats pour votre application. En général, les
diagrammes obligatoires sont les diagrammes de cas d'utilisation, de classe et de 
séquence. Vous pouvez ajouter en plus le diagramme qui vous semble pertinent :
par exemple, pour une application sur plusieurs tiers, il est intéressant de
montrer le diagramme de déploiement;
\begin{itemize}

\item Les diagrammes doivent être clairs, lisibles et bien expliqués, sans pour autant
nous submerger de détails. Des explications trop longues deviennent ennuyeuses;
\item Si un diagramme est trop grand, vous pouvez le diviser, le représenter sous
forme de plusieurs diagrammes, ou vous abstraire de certains détails. Si c'est
impossible, imprimez le sur une grande page (A3), quitte à la plier ensuite. Le
plus important est que tous les mots soient lisibles.

\end{itemize}
\subsection{Diagramme de Séquence}
Un diagramme de séquence :
\begin{itemize}
	\item Représente un scénario possible qui se déroule dans un cas d'utilisation. 
	Vous n'êtes donc pas obligés de montrer tous les cas d'exécution possibles;
	\item Représente l'interaction entre les objets : donc normalement, toutes les
	instances définies dans un diagramme de séquences doivent correspondre 
	 à des classes qui se trouvent dans le diagramme des classes;
	 \item Il existe parfois des dizaines de diagrammes de séquences possibles. Choisissez certains d'entre eux à mettre dans le rapport (2 ou 3). Privilégiez les diagrammes les plus importants (et non, l'authentification n'en fait pas partie!).
\end{itemize}
\subsection{Diagramme de Classes}
Un diagramme de classes :
\begin{itemize}
\item Doit être fidèle à l'architecture logicielle choisie. Si vous utilisez le MVC, 
alors les trois couches doivent être représentées dans le diagramme de classes grâce aux packages;
\item Les stéréotypes sont fortement conseillés. Si vous développez une
application web, n'hésitez pas à utiliser les stéréotypes de la figure \ref{fig:fig2} : 
\begin{figure}[!ht]\centering
\includegraphics[scale=0.9]{stereotypes.jpg}
\caption{Les stéréotypes}
\label{fig:fig2}
\end{figure}
\item Attention à ne pas confondre classes et tables : évitez la tentation de
mettre des id partout.

\end{itemize}
 
\section*{Conclusion}
Faire ici une petite récapitulation du chapitre, ainsi qu'une introduction du chapitre suivant.





%==============================================================================
\end{spacing}
